\documentclass{article}

\usepackage{amsmath,amssymb}
\usepackage{xurl}
\usepackage[hidelinks]{hyperref}
\setlength{\emergencystretch}{2em}

\input{command}

\title{The Inverse on the Range in the Two-Variable Operator Schwarz Inequality}
\author{}
\date{2026-08-21}

\begin{document}
\maketitle
\gapnote{open-gap}{resolved}

\section{Source statement}

Let
\begin{align}
    T&\colon\MN{d}\to\MN{d'},
    &
    E=T^*&\colon\MN{d'}\to\MN{d},
    \label{eq:wolf_ch5_two_variable_unconditional_maps}
\end{align}
where $T$ is two-positive and $E$ is its trace adjoint. Wolf's Theorem~5.3
asserts that, for every $A,B\in\MN{d'}$,
\begin{align}
    E(A^\dagger B)
    E(B^\dagger B)^{-1}
    E(B^\dagger A)
        &\leq E(A^\dagger A),
    \label{eq:wolf_ch5_two_variable_unconditional_source_inequality}
\end{align}
where the inverse is taken on the range of $E(B^\dagger B)$
\cite[Chapter~5, Theorem~5.3]{Wolf2012Quantum}.\footnote{Local source:
\path{Notes/WolfNoteTexSource/ch05_schwarz_inequalities.tex},
lines~180--190.}

Rectangular two-positivity is invariant under the trace adjoint, so $E=T^*$ is
also two-positive. The formal theorem retains the independent input and output
dimensions used by Wolf and assumes no additional range witness. It represents
the inverse on the range by the Moore--Penrose inverse.

\section{The range inverse}

Define
\begin{align}
    P&=E(A^\dagger A),
    &
    Z&=E(A^\dagger B),
    &
    R&=E(B^\dagger B).
    \label{eq:wolf_ch5_two_variable_unconditional_block_entries}
\end{align}
Two-positivity of $E$ gives the block inequality
\begin{align}
    \begin{pmatrix}
        P&Z\\
        Z^\dagger&R
    \end{pmatrix}
        &\geq 0.
    \label{eq:wolf_ch5_two_variable_unconditional_block_positive}
\end{align}
The kernel clause of the Schur-complement theorem applied to
Eq.~\ref{eq:wolf_ch5_two_variable_unconditional_block_positive} gives
\begin{align}
    \ker R&\subseteq\ker Z.
    \label{eq:wolf_ch5_two_variable_unconditional_kernel_inclusion}
\end{align}

Let $R^+$ denote the Moore--Penrose inverse of $R$, and define
\begin{align}
    s_R&=RR^+
    \label{eq:wolf_ch5_two_variable_unconditional_left_support_projection}\\
       &=R^+R.
    \label{eq:wolf_ch5_two_variable_unconditional_right_support_projection}
\end{align}
Because $R\geq 0$, $s_R$ is the orthogonal projection onto
$\operatorname{ran}R=(\ker R)^\perp$. The kernel inclusion in
Eq.~\ref{eq:wolf_ch5_two_variable_unconditional_kernel_inclusion} is
equivalent to the support-absorption identities
\begin{align}
    Zs_R&=Z,
    \label{eq:wolf_ch5_two_variable_unconditional_right_support_absorption}\\
    s_RZ^\dagger&=Z^\dagger.
    \label{eq:wolf_ch5_two_variable_unconditional_left_support_absorption}
\end{align}
Consequently,
\begin{align}
    R(R^+Z^\dagger)
        &=(RR^+)Z^\dagger
    \label{eq:wolf_ch5_two_variable_unconditional_range_witness_first}\\
        &=s_RZ^\dagger
    \label{eq:wolf_ch5_two_variable_unconditional_range_witness_second}\\
        &=Z^\dagger.
    \label{eq:wolf_ch5_two_variable_unconditional_range_witness}
\end{align}
For every vector $v$, the vector $R^+Z^\dagger v$ is therefore a preimage of
$Z^\dagger v$ under $R$. Equation
\ref{eq:wolf_ch5_two_variable_unconditional_range_witness} gives exactly the
inverse-on-the-range operation required in
Eq.~\ref{eq:wolf_ch5_two_variable_unconditional_source_inequality}. No
separate range--kernel duality argument is needed.

The formal proof establishes
Eq.~\ref{eq:wolf_ch5_two_variable_unconditional_kernel_inclusion}, then
Eqs.~\ref{eq:wolf_ch5_two_variable_unconditional_right_support_absorption}
and~\ref{eq:wolf_ch5_two_variable_unconditional_left_support_absorption}, and
finally Eq.~\ref{eq:wolf_ch5_two_variable_unconditional_range_witness}.

\section{Derivation of the inequality}

Fix a vector $v$ and define
\begin{align}
    w&=R^+Z^\dagger v.
    \label{eq:wolf_ch5_two_variable_unconditional_quadratic_witness}
\end{align}
Equation~\ref{eq:wolf_ch5_two_variable_unconditional_range_witness} gives
\begin{align}
    Rw&=Z^\dagger v.
    \label{eq:wolf_ch5_two_variable_unconditional_witness_equation}
\end{align}
The pseudoinverse-free quadratic-form consequence of
Eq.~\ref{eq:wolf_ch5_two_variable_unconditional_block_positive} then yields
\begin{align}
    v^\dagger Pv
        &\geq v^\dagger Zw
    \label{eq:wolf_ch5_two_variable_unconditional_quadratic_bound_first}\\
        &=v^\dagger ZR^+Z^\dagger v.
    \label{eq:wolf_ch5_two_variable_unconditional_quadratic_bound_second}
\end{align}
Because this inequality holds for every $v$,
\begin{align}
    ZR^+Z^\dagger&\leq P.
    \label{eq:wolf_ch5_two_variable_unconditional_schur_bound}
\end{align}
Substituting the definitions from
Eq.~\ref{eq:wolf_ch5_two_variable_unconditional_block_entries} into
Eq.~\ref{eq:wolf_ch5_two_variable_unconditional_schur_bound} gives
Eq.~\ref{eq:wolf_ch5_two_variable_unconditional_source_inequality}, with
$R^+$ realizing the inverse on the range.\footnote{Formal declaration:
\leanid{schwarz_two_variable_unconditional_of_traceAdjointMap} in the namespace
\leanid{SchwarzTwoVariable}.}

\section{Verdict}

No hypothesis or dimensional restriction remains relative to Wolf's
Theorem~5.3 \cite[Chapter~5, Theorem~5.3]{Wolf2012Quantum}. The formal theorem
records the positive map as the trace adjoint $E=T^*$ of a two-positive
reverse-direction map $T$. Rectangular two-positivity of $E$ follows from
invariance under the trace adjoint. The Moore--Penrose inverse realizes Wolf's
``inverse on the range'', and the required range inclusion follows from
support-projection absorption. The earlier proposed Euclidean
range--kernel-duality argument is unnecessary. The former gap is therefore
resolved.

\bibliographystyle{alpha}
\bibliography{references}

\end{document}
